\section{Twierdzenie Löba}\label{sec:loeb}
\begin{twierdzenie}[Löb]\label{res:13:1}
Niech $T\supseteq Q$, a ustalony predykat $\operatorname{Bew}_T$ spełnia (D1)--(D3). Dla dowolnego zdania $\varphi$:
\[
T\vdash\Box\varphi\to\varphi\quad\Longrightarrow\quad T\vdash\varphi.
\]
Do tego abstrakcyjnego wniosku nie potrzeba niesprzeczności. Założenia poprzedniej części, w szczególności standardowy predykat dla $T\supseteq\mathrm{PA}$, zapewniają wymagane warunki.
\end{twierdzenie}
\begin{proof}
Ustalmy $\varphi$ i załóżmy $T\vdash\Box\varphi\to\varphi$. Na mocy Lematu~\ref{res:8:1} weźmy punkt stały $L$ formuły $\psi(x):=(\operatorname{Bew}_T(x)\to\varphi)$, tj.
\begin{equation}\label{eq:13:1}
T\vdash L\leftrightarrow(\Box L\to\varphi).
\end{equation}
$x$ jest zmienną liczbową, dlatego w formule wejściowej występuje $\operatorname{Bew}_T(x)$, nie niepoprawny zapis $\Box x$.

\emph{1. Od $\Box L$ do $\Box\varphi$.} Korzystamy najpierw z jednego kierunku punktu stałego:
\begin{align*}
T&\vdash L\to(\Box L\to\varphi) &&\text{(z~\eqref{eq:13:1})},\\
T&\vdash\Box(L\to(\Box L\to\varphi)) &&\text{(D1)},\\
T&\vdash\Box L\to\Box(\Box L\to\varphi) &&\text{(D2)}.
\end{align*}
Osobna instancja (D2), dla zdań $\Box L$ i $\varphi$, daje $T\vdash\Box(\Box L\to\varphi)\to(\Box\Box L\to\Box\varphi)$. Jeśli teraz wewnątrz $T$ założymy $\Box L$, poprzedni ciąg daje $\Box(\Box L\to\varphi)$, a (D3) daje $\Box\Box L$. Dwa zastosowania modus ponens dają $\Box\varphi$. Usuwając pomocnicze założenie przez twierdzenie o dedukcji, otrzymujemy $T\vdash\Box L\to\Box\varphi$.

\emph{2. Od tej implikacji do $L$.} Z hipotezy twierdzenia $T\vdash\Box\varphi\to\varphi$ wynika teraz $T\vdash\Box L\to\varphi$. Odwrotny kierunek punktu stałego to $T\vdash(\Box L\to\varphi)\to L$. Jego przesłanka jest już twierdzeniem $T$, więc modus ponens daje $T\vdash L$. Nie założyliśmy ani $L$, ani $\Box L$: pomocnicze założenie z poprzedniego kroku zostało usunięte.

\emph{3. Zakończenie.} Dopiero do uzyskanego rzeczywistego dowodu $L$ stosujemy (D1), otrzymując $T\vdash\Box L$. Łącząc to z udowodnionym w kroku 2 $T\vdash\Box L\to\varphi$, dostajemy $T\vdash\varphi$.
\end{proof}

\begin{uwaga}\label{res:13:2}
Drugie twierdzenie Gödla jest szczególnym przypadkiem twierdzenia Löba dla $\varphi:=\bot$ ($0=S(0)$), gdy $T$ jest niesprzeczna: gdyby $T\vdash\operatorname{Con}_T$, czyli $T\vdash\neg\Box\bot$, to w szczególności $T\vdash\Box\bot\to\bot$, skąd twierdzenie Löba daje $T\vdash\bot$ --- sprzeczność. Zatem $T\nvdash\operatorname{Con}_T$.
\end{uwaga}

% Rozdział 14, strony 25--29.
